\section{Appendix: Matrix Algebra}

An $m\times n$ matrix has $m$ rows and $n$ columns. Addition and subtraction
require the same dimensions and operate entry by entry. The product of an
$m\times n$ matrix and an $n\times r$ matrix is an $m\times r$ matrix:
its $(i,j)$ entry is the dot product of row $i$ with column $j$.
Throughout the complex calculations, use $i^2=-1$.

\exercise{A.1}{Adding two real matrices}
\idea Add entries occupying the same position; both matrices are $2\times2$.
\[
\begin{bmatrix}-2&4\\1&5\end{bmatrix}+
\begin{bmatrix}4&-3\\1&2\end{bmatrix}
=\begin{bmatrix}-2+4&4-3\\1+1&5+2\end{bmatrix}
=\boxed{\begin{bmatrix}2&1\\2&7\end{bmatrix}}.
\]
No row is multiplied by a column when adding matrices.

\exercise{A.2}{Adding two complex matrices}
Both matrices are $2\times3$, so addition is defined. Add real parts to real
parts and imaginary parts to imaginary parts:
\begin{align*}
(1+i)+3&=4+i,& -2+(i-1)&=-3+i,&(3-i)+(-2+i)&=1,\\
4+(3-i)&=7-i,&i+(1+i)&=1+2i,&(2-i)+0&=2-i.
\end{align*}
\answer{The sum is $\displaystyle
\begin{bmatrix}4+i&-3+i&1\\7-i&1+2i&2-i\end{bmatrix}$.}

\exercise{A.3}{Subtracting complex matrices}
Subtraction is defined because both matrices are $3\times2$. Distribute each
minus sign to the entire corresponding entry:
\[
\begin{bmatrix}i&-1\\4&1\\3&-2i\end{bmatrix}-
\begin{bmatrix}3-i&4i\\2&1+i\\3&-i\end{bmatrix}
=\begin{bmatrix}i-(3-i)&-1-4i\\4-2&1-(1+i)\\3-3&-2i-(-i)\end{bmatrix}
=\boxed{\begin{bmatrix}-3+2i&-1-4i\\2&-i\\0&-i\end{bmatrix}}.
\]
In particular, subtracting $-i$ adds $i$.

\exercise{A.4}{A product of two square matrices}
\idea The inner dimensions match. Compute four row--column dot products.
\[
\begin{bmatrix}1&-1\\3&1\end{bmatrix}
\begin{bmatrix}2&4\\-1&3\end{bmatrix}
=\begin{bmatrix}1(2)+(-1)(-1)&1(4)+(-1)(3)\\
3(2)+1(-1)&3(4)+1(3)\end{bmatrix}
=\boxed{\begin{bmatrix}3&1\\5&15\end{bmatrix}}.
\]
The factors and the result are all $2\times2$.

\exercise{A.5}{A rectangular matrix product}
A $2\times2$ matrix times a $2\times3$ matrix gives a $2\times3$ result:
\[
\begin{bmatrix}3&1\\-4&2\end{bmatrix}
\begin{bmatrix}1&5&-3\\2&1&6\end{bmatrix}
=\begin{bmatrix}
3(1)+1(2)&3(5)+1(1)&3(-3)+1(6)\\
-4(1)+2(2)&-4(5)+2(1)&-4(-3)+2(6)
\end{bmatrix}
=\boxed{\begin{bmatrix}5&16&-3\\0&-18&24\end{bmatrix}}.
\]

\exercise{A.6}{Checking dimensions before multiplying}
The first matrix is $2\times2$ and the second is $3\times2$. The number of
columns of the first, namely $2$, differs from the number of rows of the
second, namely $3$. For example, the proposed first entry would require a
dot product between a two-entry row and a three-entry column.
\answer{The indicated matrix product is undefined.}

\exercise{A.7}{A complex matrix product}
Multiply rows by columns, remembering $i(3i)=3i^2=-3$:
\begin{align*}
\begin{bmatrix}i&1\\-2&1\end{bmatrix}
\begin{bmatrix}3i&1\\4&-2i\end{bmatrix}
&=\begin{bmatrix}i(3i)+1(4)&i(1)+1(-2i)\\
(-2)(3i)+1(4)&(-2)(1)+1(-2i)\end{bmatrix}\\
&=\boxed{\begin{bmatrix}1&-i\\4-6i&-2-2i\end{bmatrix}}.
\end{align*}

\exercise{A.8}{The fourth power of a real matrix}
Let $M=\begin{bmatrix}1&-1\\1&0\end{bmatrix}$. Matrix powers mean repeated
matrix multiplication, not taking powers of the individual entries.
First compute
\[
M^2=\begin{bmatrix}1-1&-1+0\\1+0&-1+0\end{bmatrix}
=\begin{bmatrix}0&-1\\1&-1\end{bmatrix}.
\]
Now square that result:
\[
M^4=(M^2)^2
=\begin{bmatrix}0&-1\\1&-1\end{bmatrix}
 \begin{bmatrix}0&-1\\1&-1\end{bmatrix}
=\boxed{\begin{bmatrix}-1&1\\-1&0\end{bmatrix}}.
\]
As a check, multiplying $M^2$ by $M$ gives $M^3=-I_2$, so $M^4=-M$.

\exercise{A.9}{The fourth power of a complex matrix}
Put $N=\begin{bmatrix}1&-i\\i&1\end{bmatrix}$. Then
\[
N^2=\begin{bmatrix}1+(-i)i&-i-i\\i+i&i(-i)+1\end{bmatrix}
=\begin{bmatrix}2&-2i\\2i&2\end{bmatrix}=2N.
\]
Consequently $N^4=(N^2)^2=(2N)^2=4N^2=8N$, and therefore
\[
\boxed{N^4=\begin{bmatrix}8&-8i\\8i&8\end{bmatrix}}.
\]
Only the scalar factors commute here; no interchange of matrix factors is needed.

\exercise{A.10}{An explicit failure of commutativity}
Take integer matrices
\[
P=\begin{bmatrix}1&0\\0&0\end{bmatrix},\qquad
Q=\begin{bmatrix}0&1\\0&0\end{bmatrix}.
\]
Their products are
\[
PQ=\begin{bmatrix}0&1\\0&0\end{bmatrix},\qquad
QP=\begin{bmatrix}0&0\\0&0\end{bmatrix}.
\]
The upper-right entries differ, so $PQ\ne QP$. Both factors belong to
$M_2(\Z)$ and both products are defined, as required.

\exercise{A.11}{Inverting a signed permutation matrix}
Let $J=\begin{bmatrix}0&1\\-1&0\end{bmatrix}$. Multiplying gives
$J^2=\begin{bmatrix}-1&0\\0&-1\end{bmatrix}=-I_2$.
Thus $J(-J)=(-J)J=I_2$. An inverse must produce the identity on both sides,
which proves
\[
\boxed{J^{-1}=-J=\begin{bmatrix}0&-1\\1&0\end{bmatrix}}.
\]

\exercise{A.12}{Inverting a diagonal matrix}
For $D=\operatorname{diag}(2,4,-1)$, choose
$E=\operatorname{diag}(1/2,1/4,-1)$. Products of diagonal matrices multiply
their corresponding diagonal entries and have zero off-diagonal entries.
Hence
\[
DE=ED=\operatorname{diag}(2\cdot\tfrac12,4\cdot\tfrac14,(-1)(-1))=I_3.
\]
This verifies directly that
\[
\boxed{D^{-1}=\begin{bmatrix}1/2&0&0\\0&1/4&0\\0&0&-1\end{bmatrix}}.
\]
The inverse is taken over $\R$ or $\C$; it need not have integer entries.

\exercise{A.13}{The determinant of a triangular matrix}
The given matrix is lower triangular: all entries above its main diagonal
are zero. Its determinant is therefore the product of its diagonal entries:
\[
\det\begin{bmatrix}3&0&0\\10&-2&0\\4&17&8\end{bmatrix}
=3(-2)(8)=\boxed{-48}.
\]
Alternatively, expanding along the first row gives
$3\det\begin{bmatrix}-2&0\\17&8\end{bmatrix}=3(-16)$, the same answer.

\exercise{A.14}{Products of invertible matrices}
\goal Given invertible $A,B\in M_n(\C)$, exhibit two-sided inverses for
$AB$ and $BA$.
\idea Reverse the order of the inverses. Matrix multiplication is associative,
which allows regrouping without changing the order of any factors.
\step{Check the inverse of $AB$.}
\[
(AB)(B^{-1}A^{-1})=A(BB^{-1})A^{-1}=AI_nA^{-1}=I_n,
\]
\[
(B^{-1}A^{-1})(AB)=B^{-1}(A^{-1}A)B=B^{-1}I_nB=I_n.
\]
Thus $(AB)^{-1}=B^{-1}A^{-1}$.
\step{Check the inverse of $BA$.}
Similarly,
\[
(BA)(A^{-1}B^{-1})=B(AA^{-1})B^{-1}=I_n,
\qquad
(A^{-1}B^{-1})(BA)=A^{-1}(B^{-1}B)A=I_n.
\]
\answer{Both products are invertible, with
$(AB)^{-1}=B^{-1}A^{-1}$ and $(BA)^{-1}=A^{-1}B^{-1}$.}
